% P19-E2 continuation: Cariani--Rips Experiment 2 and its interpretation limits.
\chapter{Possibility Premises and Conditional Interpretation}
\label{sec:p19-conditional-interpretation-experiments}

A premise can call an exception to mind without asserting that the exception
actually occurred. This distinction motivates the second experiment in
Cariani and Rips's \emph{Conditionals, Context, and the Suppression Effect}
\cite[Secs.~2.3, 5.4]{CarianiRips2017Suppression}. This chapter supplies its
evidence record and develops the logical distinctions needed to interpret it.
The first experiment is treated in
\cref{sec:p19-conditional-reasoning-suppression}. These are the article's two
new experiments; the numbered example arguments and its discussions of other
research are not additional experiments conducted in this article.

The study record follows the evidence distinctions of
\cref{sec:p19-evidence-methods}. The state calculations below are handbook
derivations, with every state explicitly stipulated. They are not observations
of participants' mental representations. The article appeared online on
1 February 2016 and in the April 2017 issue, volume~41, pages~540--589.

\section{Experiment 2: a possibility assertion as the extra premise}
\label{sec:p19-cariani-rips-experiment-two}

Let \(p\) abbreviate that it rained, \(q\) that Alicia became wet, and \(r\)
that she remembered her umbrella. In the task descriptions,
\(\mathsf{Might}(r)\) abbreviates the English possibility assertion, without
yet identifying its psychological interpretation with a particular formal
operator. The rain example's four argument structures were as follows.
Here \(p\to q\) records the conditional's position in the classical benchmark;
it does not establish that participants understood the English conditional
materially.
\[
\begin{array}{c|l|c}
\text{Argument}&\text{Premises}&\text{Conclusion}\\\hline
\text{Basic}&p\to q,\ p&q\\
\text{Contravening with might}&p\to q,\ \mathsf{Might}(r),\ p&q\\
\text{Conjunctive}&p\to q,\ (p\land r)\to\neg q,\ p&q\\
\text{Disjunctive with might}&p\to q,\ \mathsf{Might}(r),\ p&r\lor q
\end{array}
\]
The second and fourth forms replace Experiment~1's conditional about an
umbrella preventing wetness with the statement that Alicia might have
remembered her umbrella. They do not assert \(r\), and they do not explicitly
assert \(r\to\neg q\). Basic and Conjunctive arguments were unchanged from
the first experiment. All conditionals in Experiment~2 used the past tense
\cite[Sec.~2.3]{CarianiRips2017Suppression}.

\paragraph{Participants and procedure.}
The analysis included 16 introductory psychology students at Northwestern
University, participating to fulfill a course requirement. The article reports
three exclusions for having studied logic in a college class across its two
experiments together; it does not allocate these exclusions to Experiment~2.
Thus 16 is this experiment's analysis denominator, and its separate recruited
denominator cannot be obtained by adding all three exclusions to it. Age,
gender distribution and recruitment dates are not supplied in the reported
methods. Recruitment from this course does not establish a probability sample
of the adult population.

The procedure followed Experiment~1: four argument types with 12 content
variations each, giving 48 arguments per participant, displayed one at a time
on a computer in a newly randomized order for each person. The instructions
asked whether the conclusion must be true if the premises are true, using a
binary follows/does-not-follow response. This was a judgment of logical
necessity under an instruction, rather than a numerical confidence rating.
Argument type and content were repeated within participants. Randomizing the
presentation order is distinct from randomly sampling participants or randomly
assigning them to different instruction conditions; the latter manipulation
was not part of this experiment
\cite[Sec.~2.3]{CarianiRips2017Suppression}.

\paragraph{Reported responses and comparisons.}
The primary text reports 92\% endorsement for Basic arguments and 56\% for
Contravening arguments with \emph{might}. Their difference is approximately
36 percentage points, computed from the reported rounded percentages. The
article reports the overall argument-type result
\[
F(3,45)=10.57,\qquad MS_e=0.092,\qquad p<0.001.
\]
Its Tukey HSD comparisons at \(\alpha=0.05\) report higher endorsement for
Basic than for the might-Contravening form, higher endorsement for Basic than
for Conjunctive, and higher endorsement for the might-Disjunctive form than
for Conjunctive. No other pairwise differences were significant by that test
\cite[Sec.~5.4, Fig.~4]{CarianiRips2017Suppression}.

The accessible primary prose does not print exact Experiment~2 percentages
for the Conjunctive and might-Disjunctive forms; it directs the reader to
Figure~4. Their figure-only numerical values are not transcribed here because
the figure asset was unavailable in the acquired publisher text. The same
text gives an image-only effect-size symbol followed by \(0.413\); without
the symbol's readable source image, that number is not assigned a statistic
name here. These source-access limits do not change the stated comparisons.

The design supplies a nominal \(16\times12=192\) response opportunities per
argument type and \(16\times48=768\) across types. These products are design
denominators, not published raw response counts or numbers of independent
participants. The paper's reported rounded rates do not identify each
participant's twelve responses, missing-response pattern, or covariance
between conditions. The error mean square is the reported analysis quantity;
it is not by itself a standard deviation of binary responses or a confidence
interval for the 36-point contrast. The acquired text supplies no participant
records or confidence interval for that contrast. A nonsignificant comparison
does not establish equality or equivalence of its population rates.

\paragraph{What the comparison establishes.}
In this sample and task, a possibility premise accompanied lower endorsement
of the unchanged categorical conclusion than the Basic form. An explicit
second conditional was therefore not required for this observed suppression.
The result remains compatible with different accounts of why the possibility
premise mattered. It did not directly measure which alternatives participants
considered, whether they revised the first conditional, or how they selected
their final response. Most notably, the observed 56\% is not a finding that
every participant suppressed the inference, nor is it the proportion of
participants who always did so.

Experiment~1's corresponding conditional-premise rate was 35\%, but comparing
35\% and 56\% across the two experiments does not isolate the causal effect
of replacing the conditional with \emph{might}. The groups differed, and the
unchanged Basic and Conjunctive forms also produced less extreme results in
Experiment~2. The authors explicitly caution against giving the between-study
differences too much weight \cite[Sec.~5.4]{CarianiRips2017Suppression}.

\section{A possibility assertion, a fixed interpretation and a changed state}
\label{sec:p19-might-and-fixed-context}

\begin{hbdefinition}{def:p19-possibility-state}{Finite states and compatibility}
For the following derivations, \(W=\{0,1\}^{\{p,q,r\}}\) is the set of
valuations of three atoms, extended to nonmodal propositional formulas by
the classical clauses of Section~\ref{sec:prop-syntax-semantics-02},
with \(0\) and \(1\) coding false and true, respectively. Write
\([A]=\{v\in W:v(A)=1\}\). A state \(I\) is a nonempty subset of \(W\).
It \emph{accepts} \(A\) when \(I\subseteq[A]\), and it is
\emph{compatible} with \(A\), written \(\operatorname{Poss}_I(A)\), when
\(I\cap[A]\ne\varnothing\). These are universal and existential conditions,
respectively. Compatibility is a minimal finite interpretation of a
possibility claim; it is not a complete linguistic semantics for every use
of \emph{might}.

For nonmodal \(A,B\), let \(C_I(A,B)\) mean
\(I\cap[A]\subseteq[B]\), the restricted strict-conditional condition at
the fixed state \(I\). Define categorical restriction
\(E_A(S)=S\cap[A]\) and conditional exclusion
\(E_{A,B}(S)=S\setminus([A]\cap[\neg B])\) for every subset
\(S\subseteq W\), including the empty set.
An empty result represents incompatible accepted constraints, not a
subsequent nonempty state. These definitions leave an expansion operation
unspecified: an expansion requires a choice of newly considered valuations.
\end{hbdefinition}

\begin{hblemma}{thm:p19-fixed-state-mp}{proved=annex:p19.fixed-state-mp}{Fixed-state MP and compatible exceptions}
If a state \(I\) satisfies \(C_I(p,q)\) and accepts \(p\), it accepts \(q\).
Adding \(\operatorname{Poss}_I(r)\) to these conditions preserves this
conclusion. Moreover,
\[
 E_p(E_{p,q}(I))=E_{p,q}(E_p(I))=I\cap[p]\cap[q].
\]
If the result is nonempty it accepts \(q\), in either processing order.
Finally, \(C_I(p,q)\) and \(C_I(p\land r,\neg q)\) together imply
\(I\cap[p]\cap[r]=\varnothing\). They are incompatible with the further
condition \(\operatorname{Poss}_I(p\land r)\), but can both hold when that
compatibility condition is absent.
There can also be a valuation \(w\notin I\) with \(C_I(p,q)\) true
but the material conditional \(p\to q\) false at \(w\).
\end{hblemma}

The last distinction matters when interpreting apparently conflicting
conditionals. The pair ``rain implies wet'' and ``rain with an umbrella
implies dry'' can be vacuously compatible on a state containing no rainy
umbrella possibility. A requirement that the second antecedent be possible
changes this conclusion. A source explanation invoking such a requirement
must retain it \cite[Sec.~4.1, notes~3--4]{CarianiRips2017Suppression}.
Likewise, a strict conditional evaluated over a state that
does not contain the actual world need not yield the actual-world material
conditional. Our fixed-state result concerns acceptance throughout the state,
not an unproved claim about an actual world outside it.

\begin{hbexample}{def:p19-might-expansion-example}{The same possibility claim allows different expansions}
List a valuation as \((p,q,r)\). Begin with
\(I_0=\{(0,0,0),(1,1,0)\}\), where \(C_{I_0}(p,q)\) holds and
\(r\) is not possible. Two choices of expansion are
\[
 J_{\rm wet}=I_0\cup\{(1,1,1)\},\qquad
 J_{\rm dry}=I_0\cup\{(1,0,1)\}.
\]
Both accommodate a possibility that \(r\) holds. After accepting \(p\),
their respective states are
\[
 K_{\rm wet}=\{(1,1,0),(1,1,1)\},\qquad
 K_{\rm dry}=\{(1,1,0),(1,0,1)\}.
\]
These tables are constructed examples, with no empirical weights attached.
\end{hbexample}

\begin{hblemma}{thm:p19-might-expansion}{proved=annex:p19.might-expansion}{Possibility alone does not determine suppression}
For each final state \(K\) in \cref{def:p19-might-expansion-example},
\(\operatorname{Poss}_K(r)\) holds; both accept \(p\) and \(r\lor q\) and fail to
accept \(r\). The wet state accepts \(q\); the dry state accepts neither
\(q\) nor \(\neg q\). The original restricted conditional remains true
on the wet state and is false on the dry state. Thus the requirement that
\(r\) be possible alone neither asserts \(r\) nor determines whether
\(q\) is accepted after the expansion. On the dry state, adding the
categorical information \(\neg r\) leaves just \((1,1,0)\), which accepts
\(q\).
\end{hblemma}

The source's account selects a dry possibility through the conditional's
content and background expectations about umbrellas. Experiment~2 does not
supply the disabling conditional explicitly, so that selection must be
explained by the interpretation account. It cannot be smuggled in as the
literal meaning of \(\mathsf{Might}(r)\). The dry construction illustrates
how the weaker disjunction can survive when \(q\) alone does not. It also
exposes the precise change: the first conditional is no longer satisfied in
the expanded state. Classical MP with fixed premises and meanings remains
valid by \cref{thm:p19-conditional-benchmark}; the calculation has changed the
state used to interpret a premise.

\section{The source's explanatory commitments and remaining questions}
\label{sec:p19-conditional-context-explanations}

Cariani and Rips distinguish propositions taken for granted in a conversation,
the context set of worlds compatible with them, and the salient state used
to evaluate a conditional. The salient state must be a subset of the context
set; the two are not identified. Their account initially favors typical
possibilities and updates the salient state while premises are interpreted.
Categorical assertions restrict possibilities. Bare conditionals exclude
antecedent-true, consequent-false possibilities. An antecedent's required
compatibility, or an epistemic possibility assertion, can prompt expansion
\cite[Secs.~3--5.4]{CarianiRips2017Suppression}.

Its Composite Possibility assumption further favors combining a newly salient
possibility with each side of another salient distinction, unless contextual
reasons rule a combination out. Thus thinking about an umbrella can prompt
both rainy and nonrainy umbrella cases. This is an additional constraint on
which cases are considered, not a consequence of existential compatibility
alone. Background knowledge and typicality still determine which combinations
are admitted. The two expansions above deliberately hold the compatibility
requirement constant while exposing this extra choice.

This is a proposal about how semantic and pragmatic processing interact.
It does not require literal enumeration of worlds or a particular mental
diagram. Nor does the finite construction above supply an independently
measured ranking of typicality, a probability distribution over expansions,
or a model predicting the numerical 92\% and 56\% rates. Those would be
further psychological claims needing their own definitions and evidence.

\paragraph{Reversed premise order.}
The article's Negative Possibility principle asks that the state be compatible
with the negation of a conditional's antecedent. The authors identify an
objection to treating this as an unconditional update rule. After accepting \(p\), a later
conditional about \(p\) could mechanically restore \(\neg p\) possibilities,
including \(\neg p\land\neg q\). That would undermine the conclusion of
reversed MP. The authors instead allow the conditional to signal that
acceptance of \(q\) depends on \(p\), without reinstating a possibility
just excluded by the categorical premise. They acknowledge more than one
pragmatic response to the same apparent violation
\cite[Sec.~5.3]{CarianiRips2017Suppression}.
The commutation calculation in \cref{thm:p19-fixed-state-mp} identifies the
stable mathematical case. A theory adding expansion needs a further rule
for when expansion is appropriate. Experiment~2 did not manipulate reversed
premise order and cannot be counted as a direct test of that repair.

\paragraph{Conditional perfection and other inference forms.}
The article proposes that some speakers interpret a conditional more strongly,
as exhausting the sufficient conditions for its consequent. Such
\emph{conditional perfection} can help explain acceptance of AC or DA;
an explicitly supplied alternative sufficient condition can challenge the
strengthening. Its distribution is presented as a theoretical sketch,
requiring further evidence and a more precise account of sufficiency
\cite[Sec.~5.6]{CarianiRips2017Suppression}.
The proposed pragmatic licenses are an expectation that the speaker is
exhausting sufficient conditions and the absence of another salient sufficient
condition. Mentioning a pond supplies a competing way to become wet and can
defeat that interpretation. Mentioning whether the umbrella was forgotten
need not introduce an independent source of wetness. These are claims about
what interlocutors take as sufficient; merely displaying two conditionals
does not prove that either licensing condition is psychologically present.
Experiment~2's four forms are variants of MP, so it supplies no new AC or
DA sample. The benchmark in \cref{thm:p19-conditional-benchmark} applies to
the unstrengthened material conditional. It does not classify reasoning
from an explicitly assumed biconditional as the same argument.

\paragraph{Comparing explanations.}
The article contrasts its contextual account with mental-model modulation,
default or abnormality representations, conversational implicatures, and
probability-based accounts of response. Its objections concern how each
account selects relevant possibilities or probabilities and predicts the
different conclusion forms. These are the authors' arguments about rival
accounts, not additional primary experiments reproduced here. The article
also leaves unification of categorical acceptance with conditional exclusion
open, and describes its broader argument about semantic theory as requiring
refinement \cite[Secs.~4.2, 6.1--6.2]{CarianiRips2017Suppression}.

In the mental-model comparison, adding umbrella information changes the
represented possibilities so that rain together with forgetting the umbrella
becomes the operative condition. Cariani and Rips share the appeal to
possibilities and context, but propose a fixed conditional clause whose state
parameter changes. Their criticism of unrestricted modulation is that it
needs constraints on when possibilities are added, deleted or reinterpreted.
Their own state-selection principles must meet that same explanatory demand.
Experiment~2's aggregate responses do not decide whether the representation
was a diagram, mental sentences, or another format
\cite[Sec.~6.1.1]{CarianiRips2017Suppression}.

In the default-account comparison, an antecedent produces its consequence
only when no relevant abnormality intervenes. Umbrella information can mark
an abnormality for rain's producing wetness. A closed-world assumption then
treats an appropriate list of abnormalities as exhaustive. The source's
objection concerns the justification of this closure in ordinary discourse:
raincoats, shelter and other unmentioned factors could also matter. A database
whose records are stipulated exhaustive supplies a different justification
from a conversational inference that an unmentioned exception is absent.
This comparison challenges the specification of relevance and exhaustiveness;
it is not a proof that every default formalism predicts the wrong response
\cite[Sec.~6.1.2]{CarianiRips2017Suppression}.

In the implicature comparison, an added premise is said to undermine an
earlier implication that rain is sufficient for wetness. The source asks
whether this sufficiency belongs to the conditional's content or to a
cancelable conversational implication. It also asks how merely canceling
sufficiency accounts for endorsement of the weaker disjunctive conclusion.
These objections require an account of what information remains available
after the change. Our dry table gives one explicit remaining-information
pattern; it does not establish that participants used it
\cite[Sec.~6.1.3]{CarianiRips2017Suppression}.

The probability-account comparison distinguishes predicting responses from
the unconditional probability of the conclusion and predicting them from
its probability given the categorical premise. Neither quantity, without
an account of how it changes when the possibility premise is read, supplies
a full update mechanism. The authors accept that probability judgments can
matter when premises themselves are uncertain, while questioning whether
they explain every necessity judgment. Their argument does not turn the
observed frequency of endorsing an inference into a participant's subjective
probability that Alicia became wet. Those are different measured or proposed
quantities \cite[Sec.~6.1.4]{CarianiRips2017Suppression}.

The comparison yields a useful research question: which independently
specified interpretation rule predicts a participant's response across
Basic, possibility, Conjunctive and Disjunctive tasks? Agreement with the
direction of one contrast does not choose a unique rule. Measuring or
manipulating the assumed background, keeping the response criterion explicit,
and assessing predictions on new items would supply different evidence from
fitting a story after seeing the aggregate percentages. This is a proposed
evidence requirement, not a completed follow-up experiment.

\paragraph{Instructions and generalization.}
The paper's discussion of implicit suppression proposes task instructions and
response scales as possible contributors to differences between earlier
studies. It also acknowledges earlier reports of implicit suppression under
deduction instructions. Those cited studies have not been added to this
chapter's primary-study sample, and their reported differences do not become
a randomized instruction experiment by being compared in a discussion
\cite[Sec.~6.1.5]{CarianiRips2017Suppression}.
The warranted record is an effect in the stated student sample under the
stated necessity judgment. Transfer to other languages, populations, expert
reasoning, spontaneous conversation or repeated training requires further
evidence. The result informs how an English reasoning task is interpreted;
it supplies no countermodel to a proved deductive rule.
